% © 2026 Ing. David Jaroš — CC BY-NC-ND 4.0
%
% This work is licensed under a Creative Commons Attribution-NonCommercial-NoDerivatives
% 4.0 International License (CC BY-NC-ND 4.0).
%
% License History: Earlier drafts (up to v0.3) were released under CC BY 4.0.
% From v0.4 onward, all material is released under CC BY-NC-ND 4.0 to protect
% the integrity of the theoretical work during ongoing academic development.

\documentclass[12pt,a4paper]{article}

\usepackage{amsmath,amssymb,amsthm}
\usepackage{mathtools}
\usepackage{geometry}
\usepackage{booktabs}
\usepackage{hyperref}
\usepackage{xcolor}

\geometry{left=2.5cm,right=2.5cm,top=2.5cm,bottom=2.5cm}

\definecolor{provengreen}{rgb}{0.0,0.50,0.0}
\definecolor{openred}{rgb}{0.70,0.0,0.0}
\definecolor{semiblue}{rgb}{0.0,0.0,0.65}

\newtheorem{theorem}{Theorem}[section]
\newtheorem{lemma}[theorem]{Lemma}
\newtheorem{proposition}[theorem]{Proposition}
\newtheorem{corollary}[theorem]{Corollary}
\theoremstyle{definition}
\newtheorem{definition}[theorem]{Definition}
\theoremstyle{remark}
\newtheorem{remark}[theorem]{Remark}

\newcommand{\C}{\mathbb{C}}
\newcommand{\R}{\mathbb{R}}
\renewcommand{\H}{\mathbb{H}}
\newcommand{\Neff}{N_{\mathrm{eff}}}
\newcommand{\Tr}{\mathrm{Tr}}

\title{%
  $N_\mathrm{eff}$ Derivation, Step 3: Final Result and Derivation Status\\
  \large $N_{\mathrm{eff}}$ from UBT Action $S[\Theta]$ --- $\alpha$/B-coefficient Track}
\author{David Jaroš}
\date{2026-03-05}

\begin{document}
\maketitle

\begin{abstract}
This document states the final result of the $\Neff$ derivation chain
(Steps~1--3) and updates the derivation status for $\Neff = 12$ and
$B_0 = 8\pi$ in the UBT $\alpha$ derivation.
The key finding is \textbf{Result A}: $\Neff = 12$ is proved from
$\C \otimes \H$ algebra alone; $B_0 = 8\pi \approx 25.13$ is a
zero-free-parameter one-loop result.
The discrepancy with the phenomenologically required $B \approx 46.3$
is precisely characterised as Open Problem A, \emph{unchanged} by this derivation.
\end{abstract}

\tableofcontents

%──────────────────────────────────────────────────────────────────────────────
\section{Summary of Steps 1--2}
\label{sec:summary}
%──────────────────────────────────────────────────────────────────────────────

\begin{center}
\begin{tabular}{llll}
  \toprule
  \textbf{Step} & \textbf{File} & \textbf{Result} & \textbf{Status} \\
  \midrule
  1 & \texttt{step1\_mode\_decomposition.tex}
    & $N_\mathrm{phases} = 3$ from $\H$; $\Neff = 12 = 3 \times 2 \times 2$
    & \textcolor{provengreen}{\textbf{Proved}} \\
  2 & \texttt{step2\_vacuum\_polarization.tex}
    & $B_0 = 2\pi\Neff/3 = 8\pi \approx 25.13$ from $S_\mathrm{kin}[\Theta]$
    & \textcolor{provengreen}{\textbf{Proved}} \\
  3 & This document
    & Result classification; DERIVATION\_INDEX update
    & \textcolor{provengreen}{\textbf{Complete}} \\
  \bottomrule
\end{tabular}
\end{center}

%──────────────────────────────────────────────────────────────────────────────
\section{The Three Possible Results}
\label{sec:results}
%──────────────────────────────────────────────────────────────────────────────

From the problem statement (Track CORE, $\alpha$/B-coefficient, 2026-03-05),
the three possible outcomes were:

\begin{enumerate}
  \item[\textbf{A.}] $\Neff = 12$ proved from $\C \otimes \H$ algebra alone.
    $\Rightarrow$ $B_0 = 25.1$ is a zero-free-parameter result.
    Update DERIVATION\_INDEX.
  \item[\textbf{B.}] $\Neff = 4$ from $\C \otimes \H$ (EW only),
    $\Neff = 12$ needs SU(3).
    $\Rightarrow$ Document precisely: ``$B_0 = 8.9$ from $\C \otimes \H$ alone;
    full $B_0 = 25.1$ requires SU(3) completion (Track B).''
  \item[\textbf{C.}] $\Neff = $ other value.
    $\Rightarrow$ Document with full calculation.
\end{enumerate}

%──────────────────────────────────────────────────────────────────────────────
\section{Result A: $\Neff = 12$ from $\C \otimes \H$ Alone}
\label{sec:result_A}
%──────────────────────────────────────────────────────────────────────────────

\begin{center}
  \colorbox{provengreen!15}{%
    \parbox{0.88\textwidth}{%
      \large\textbf{Result A confirmed:}
      $\Neff = 12$ is proved from $\C \otimes \H$ algebra alone.\\[4pt]
      $B_0 = 2\pi\Neff/3 = 8\pi \approx 25.133$ is a zero-free-parameter
      one-loop result.\\[4pt]
      SU(3)/octonions are \emph{not} required for this count.}}
\end{center}

\subsection{The key algebraic argument}

The factor $N_\mathrm{phases} = 3$ comes from the \emph{three imaginary
directions of $\H$} (units $\mathbf{i}, \mathbf{j}, \mathbf{k}$), which
compactify to three independent circles.  This is a property of the
quaternion algebra $\H$, present already in $\C \otimes \H$, long before
any octonionic extension.

\begin{theorem}[Main result]
\label{thm:main}
Let $\Theta \in \C \otimes_\R \H \cong \mathrm{Mat}(2,\C)$ be the UBT field
with kinetic action $S_\mathrm{kin}[\Theta] = \int d^4x\,d\psi\;
\mathrm{Sc}[(D_\mu\Theta)^\dagger(D^\mu\Theta)]$.
Under canonical compactification $\psi,\chi,\xi \sim 2\pi$ and the
$U(1)_\mathrm{EM}$ gauge symmetry, the number of independent complex charged
degrees of freedom contributing to the one-loop vacuum polarization is
\begin{equation}
  \Neff = N_\mathrm{phases} \times N_\mathrm{helicity} \times N_\mathrm{charge}
  = 3 \times 2 \times 2 = 12,
  \label{eq:Neff_main}
\end{equation}
and the one-loop $\beta$-function coefficient is
\begin{equation}
  B_0 = \frac{2\pi\Neff}{3} = 8\pi \approx 25.133.
  \label{eq:B0_main}
\end{equation}
Both values follow from the algebra of $\C \otimes \H$ alone, with zero
free parameters.
\end{theorem}

\begin{proof}
  See \texttt{step1\_mode\_decomposition.tex} (Theorem~1.4 for $\Neff = 12$)
  and \texttt{step2\_vacuum\_polarization.tex} (Theorem~3.1 for $B_0$).
\end{proof}

\subsection{Why not Result B?}

Result B would apply if the factor of 3 required SU(3) (colour).
But $N_\mathrm{phases} = 3$ comes from $\dim_\R \mathrm{Im}(\H) = 3$,
not from the number of colours.  The factor $3 = \dim(\mathrm{Im}(\H))$
is a property of the quaternion algebra alone.

Coincidentally, $N_\mathrm{colours} = 3$ (from SU(3)$_c$) has the same
numerical value, which is a structural feature of UBT: the three compact
quaternionic time directions \emph{mirror} the three quark colours without
being identified with them.  The derivation does not use SU(3).

\subsection{Implications for $B_0 = 25.1$}

The value $B_0 = 8\pi \approx 25.133$ is now a \textbf{proved} quantity:
\begin{itemize}
  \item No free parameters.
  \item QED limit: $\Neff = 1$ gives $B_0 = 2\pi/3 \approx 2.09$. $\checkmark$
  \item SM interpretation: $\Neff = 12$ corresponds to 12 SM gauge bosons
    ($8 + 3 + 1$), but this identification is a \emph{consistency check},
    not a required input.
\end{itemize}

%──────────────────────────────────────────────────────────────────────────────
\section{Open Problem A: $B_\mathrm{phenom} \approx 46.3 \neq B_0$}
\label{sec:open_problem}
%──────────────────────────────────────────────────────────────────────────────

The derivation of $\Neff = 12$ and $B_0 = 8\pi$ does \emph{not} resolve
the discrepancy with the phenomenologically required value.

\begin{center}
  \colorbox{openred!15}{%
    \parbox{0.88\textwidth}{%
      \textbf{Open Problem A (unchanged):}\\
      $B_0 = 8\pi \approx 25.13$ (derived, this work).\\
      $B_\mathrm{phenom} \approx 46.3$ (required for $n^* = 137$,
      from effective potential analysis).\\[4pt]
      Ratio: $B_\mathrm{phenom}/B_0 \approx 1.844$.\\[4pt]
      This ratio is \emph{not} a two-loop QED correction
      (which is $O(\alpha/\pi) \approx 0.002$).\\
      The discrepancy remains an open problem.
      See \texttt{STATUS\_ALPHA.md §9} and
      \texttt{appendix\_ALPHA\_one\_loop\_biquat.tex §B.3}.}}
\end{center}

The precise statement is:
\begin{itemize}
  \item $B_0 = 2\pi\Neff/3 = 8\pi$ is the \textbf{one-loop photon vacuum
    polarization} coefficient, derived rigorously from $S_\mathrm{kin}[\Theta]$.
  \item The coefficient $B \approx 46.3$ appears in the effective potential
    $V_\mathrm{eff}(n) = An^2 - Bn\ln n$ governing the topological prime
    selection.  This $B$ involves a \emph{different} mechanism (geometric
    compactification enhancement, or higher-loop effects) not captured by
    the one-loop photon self-energy alone.
  \item Until this gap is resolved, the relation
    $B_0 \approx 25.13$ is proved; $B \approx 46.3$ remains semi-empirical.
\end{itemize}

%──────────────────────────────────────────────────────────────────────────────
\section{DERIVATION\_INDEX Update}
\label{sec:index_update}
%──────────────────────────────────────────────────────────────────────────────

The following changes to \texttt{DERIVATION\_INDEX.md} are mandated by
this derivation chain:

\begin{center}
\begin{tabular}{p{4.5cm}p{3.5cm}p{4cm}p{2.5cm}}
  \toprule
  \textbf{Entry} & \textbf{Old status} & \textbf{New status} & \textbf{Reference} \\
  \midrule
  $\Neff = 12$ from SM gauge group
  & Semi-empirical [L0]
  & \textcolor{provengreen}{\textbf{Proved [L0]}}
  & Step~1 (this chain) \\[4pt]
  $B_0 = 25.1$ (one-loop baseline)
  & Proved [L1]
  & \textcolor{provengreen}{\textbf{Proved [L1]}} (unchanged)
  & Step~2 (this chain) \\[4pt]
  $B_\mathrm{phenom} \approx 46.3$ required for $n^*=137$
  & (implicit)
  & \textcolor{openred}{\textbf{Open Hard Problem}}
  & STATUS\_ALPHA.md §9 \\
  \bottomrule
\end{tabular}
\end{center}

The key change: the entry ``$\Neff = 12$ from SM gauge group''
is upgraded from \textbf{Semi-empirical} to \textbf{Proved [L0]},
with the derivation chain
\texttt{N\_eff\_derivation/step1\_mode\_decomposition.tex} $\to$
\texttt{step2\_vacuum\_polarization.tex} $\to$
\texttt{step3\_N\_eff\_result.tex}
as the canonical reference.

\paragraph{Precise statement for DERIVATION\_INDEX.}
\begin{quote}
\textit{$\Neff = 12$ is proved from $\C \otimes \H$ algebra alone as
$3 \times 2 \times 2$ (three imaginary quaternion phases $\times$ two
helicities $\times$ particle/antiparticle).  The factor 3 comes from
$\dim \mathrm{Im}(\H) = 3$, not from SU(3).
Reference: \texttt{consolidation\_project/N\_eff\_derivation/step1\_mode\_decomposition.tex},
Theorem~1.4.}
\end{quote}

%──────────────────────────────────────────────────────────────────────────────
\section{Complete Result Table}
\label{sec:full_table}
%──────────────────────────────────────────────────────────────────────────────

\begin{center}
\begin{tabular}{lllp{4cm}}
  \toprule
  \textbf{Quantity} & \textbf{Value} & \textbf{Status} & \textbf{Notes} \\
  \midrule
  $N_\mathrm{phases}$ & 3 & \textcolor{provengreen}{Proved} & $\dim\mathrm{Im}(\H)=3$ \\
  $N_\mathrm{helicity}$ & 2 & \textcolor{provengreen}{Proved} & Transverse photon polarizations \\
  $N_\mathrm{charge}$ & 2 & \textcolor{provengreen}{Proved} & CPT / particle+antiparticle \\
  $\Neff$ & 12 & \textcolor{provengreen}{\textbf{Proved [L0]}} & $3 \times 2 \times 2$, zero free params \\
  $B_0 = 2\pi\Neff/3$ & $8\pi \approx 25.13$ & \textcolor{provengreen}{\textbf{Proved [L1]}} & One-loop, zero free params \\
  $B_\mathrm{phenom}$ & $\approx 46.3$ & \textcolor{openred}{Open Problem A} & Needed for $n^*=137$ \\
  $B_\mathrm{base} = \Neff^{3/2}$ & $12^{3/2} \approx 41.57$ & \textcolor{openred}{Open Problem A} & No derivation exists \\
  \midrule
  QED check: $\Neff=1$ & $B_0 = 2\pi/3$ & \textcolor{provengreen}{Verified} & $\checkmark$ \\
  \bottomrule
\end{tabular}
\end{center}

\end{document}
